% GENERATED FILE. Edit canonical lessons, then run npm run generate:books.
% canonical-source-hash: fnv1a64:01f270a5ca5e3ff4
% canonical-lessons: ES-C645-pinchazo, ES-C645-averiarse, ES-C645-mecanico, ES-C645-neumatico, ES-C645-pinchar

\chapter{Una avería}
\label{ch:es-una-averia}
\hlchaptermodality{645}

\begin{chapteropening}
\textbf{By the end of this chapter:} \emph{I can explain that the car has broken down or has a flat tyre and ask for a mechanic.}

Complete the last lesson of chapter 645: I can explain that the car has broken down or has a flat tyre and ask for a mechanic.
\end{chapteropening}

\section[el pinchazo]{el pinchazo — the puncture, the flat tyre}

\label{lesson:ES-C645-pinchazo}



\begin{quote}
Before the new one: say the Spanish for the belt, then the Spanish for to fasten.
\end{quote}

\subsection*{You'll want to know: el pinchazo}
\textbf{el pinchazo} — \textquotedblleft{}the puncture, the flat tyre\textquotedblright{}.

\begin{quote}
Tuvimos un pinchazo cerca del pueblo.
\end{quote}

\begin{quote}
\emph{We had a puncture near the village.}
\end{quote}

In Mexico people say \emph{la ponchadura}.

\begin{cousinweb}
\emph{pinchazo} is built on \emph{pinchar} \textquotedblleft{}to prick\textquotedblright{}. The origin of \emph{pinchar} is not settled; it is often linked to Latin \emph{punctum} \textquotedblleft{}a prick\textquotedblright{}.
\end{cousinweb}

\begin{grammarlens}[title={What you've built}]
One more word for telling, explaining and describing at B1.
\end{grammarlens}

\subsection*{Your turn}
\begin{itemize}
\raggedright
  \item \emph{Say it:} \emph{el pinchazo}
  \item \emph{Say it:} \emph{el pinchazo}, once more
  \item \emph{Say it:} a sentence of your own with \emph{el pinchazo}
  \item \emph{Recall:} say the Spanish for the belt, then the Spanish for to fasten, then say \emph{el pinchazo} again
\end{itemize}

\begin{tcolorbox}[breakable,colback=teal!4,colframe=teal!35!black,title={Before you move on}]
What does el pinchazo mean? (\textquotedblleft{}The puncture, the flat tyre\textquotedblright{}.) Say it once more.
\end{tcolorbox}
\section[averiarse]{averiarse — to break down}

\label{lesson:ES-C645-averiarse}



\begin{quote}
Before the new one: say the Spanish for to fasten, then the Spanish for the puncture.
\end{quote}

\subsection*{You'll want to know: averiarse}
\textbf{averiarse} — \textquotedblleft{}to break down\textquotedblright{}.

\begin{quote}
El autobús se averió en la autopista.
\end{quote}

\begin{quote}
\emph{The bus broke down on the motorway.}
\end{quote}

Machines \emph{se averían}; a person \emph{se pone enfermo}.

\begin{cousinweb}
\emph{averiarse} is built on \emph{avería} \textquotedblleft{}breakdown\textquotedblright{}, a word that came through Catalan or Italian and is usually traced to Arabic \emph{ʿawār} \textquotedblleft{}a defect\textquotedblright{}.
\end{cousinweb}

\begin{grammarlens}[title={What you've built}]
One more word for telling, explaining and describing at B1.
\end{grammarlens}

\subsection*{Your turn}
\begin{itemize}
\raggedright
  \item \emph{Say it:} \emph{averiarse}
  \item \emph{Say it:} \emph{averiarse}, once more
  \item \emph{Say it:} a sentence of your own with \emph{averiarse}
  \item \emph{Recall:} say the Spanish for to fasten, then the Spanish for the puncture, then say \emph{averiarse} again
\end{itemize}

\begin{tcolorbox}[breakable,colback=teal!4,colframe=teal!35!black,title={Before you move on}]
What does averiarse mean? (\textquotedblleft{}To break down\textquotedblright{}.) Say it once more.
\end{tcolorbox}
\section[el mecánico]{el mecánico — the mechanic}

\label{lesson:ES-C645-mecanico}



\begin{quote}
Before the new one: say the Spanish for the puncture, then the Spanish for to break down.
\end{quote}

\subsection*{You'll want to know: el mecánico}
\textbf{el mecánico} — \textquotedblleft{}the mechanic\textquotedblright{}.

\begin{quote}
El mecánico viene en media hora.
\end{quote}

\begin{quote}
\emph{The mechanic is coming in half an hour.}
\end{quote}

As an adjective it means mechanical. A woman is \emph{la mecánica}.

\begin{cousinweb}
\emph{mecánico} is a learned borrowing from Latin \emph{mēchanicus} \textquotedblleft{}mechanical\textquotedblright{} (Lewis \& Short), which comes from Greek \emph{mēkhanikós}, built on \emph{mēkhan\'{\={e}}} \textquotedblleft{}machine\textquotedblright{} (EtymDB).
\end{cousinweb}

\begin{grammarlens}[title={What you've built}]
One more word for telling, explaining and describing at B1.
\end{grammarlens}

\subsection*{Your turn}
\begin{itemize}
\raggedright
  \item \emph{Say it:} \emph{el mecánico}
  \item \emph{Say it:} \emph{el mecánico}, once more
  \item \emph{Say it:} a sentence of your own with \emph{el mecánico}
  \item \emph{Recall:} say the Spanish for the puncture, then the Spanish for to break down, then say \emph{el mecánico} again
\end{itemize}

\begin{tcolorbox}[breakable,colback=teal!4,colframe=teal!35!black,title={Before you move on}]
What does el mecánico mean? (\textquotedblleft{}The mechanic\textquotedblright{}.) Say it once more.
\end{tcolorbox}
\section[el neumático]{el neumático — the tyre}

\label{lesson:ES-C645-neumatico}



\begin{quote}
Before the new one: say the Spanish for to break down, then the Spanish for the mechanic.
\end{quote}

\subsection*{You'll want to know: el neumático}
\textbf{el neumático} — \textquotedblleft{}the tyre\textquotedblright{}.

\begin{quote}
Hay que cambiar el neumático.
\end{quote}

\begin{quote}
\emph{The tyre has to be changed.}
\end{quote}

Spain says \emph{el neumático} or \emph{la rueda}; much of Latin America says \emph{la llanta}.

\begin{cousinweb}
\emph{neumático} is a learned borrowing from Latin \emph{pneumaticus} \textquotedblleft{}of air or wind\textquotedblright{} (Lewis \& Short), from Greek \emph{pneumatikós}, built on \emph{pneûma} \textquotedblleft{}air, breath\textquotedblright{} (EtymDB). A tyre is full of air.
\end{cousinweb}

\begin{grammarlens}[title={What you've built}]
One more word for telling, explaining and describing at B1.
\end{grammarlens}

\subsection*{Your turn}
\begin{itemize}
\raggedright
  \item \emph{Say it:} \emph{el neumático}
  \item \emph{Say it:} \emph{el neumático}, once more
  \item \emph{Say it:} a sentence of your own with \emph{el neumático}
  \item \emph{Recall:} say the Spanish for to break down, then the Spanish for the mechanic, then say \emph{el neumático} again
\end{itemize}

\begin{tcolorbox}[breakable,colback=teal!4,colframe=teal!35!black,title={Before you move on}]
What does el neumático mean? (\textquotedblleft{}The tyre\textquotedblright{}.) Say it once more.
\end{tcolorbox}
\section[pinchar]{pinchar — to get a puncture; to prick; to click}

\label{lesson:ES-C645-pinchar}



\begin{quote}
Before the new one: say the Spanish for the mechanic, then the Spanish for the tyre.
\end{quote}

\subsection*{You'll want to know: pinchar}
\textbf{pinchar} — \textquotedblleft{}to get a puncture; to prick; to click\textquotedblright{}.

\begin{quote}
Pinchamos una rueda en el camino.
\end{quote}

\begin{quote}
\emph{We got a puncture on the way.}
\end{quote}

In Spain \emph{pinchar} is also to click on a screen. Mexico says \emph{ponchar} for a puncture.

\begin{cousinweb}
The origin of \emph{pinchar} is not settled. It is often linked to Latin \emph{punctum} \textquotedblleft{}a prick, a point\textquotedblright{}, but the path is unclear.
\end{cousinweb}

\begin{grammarlens}[title={What you've built}]
One more word for telling, explaining and describing at B1.
\end{grammarlens}

\subsection*{Your turn}
\begin{itemize}
\raggedright
  \item \emph{Say it:} \emph{pinchar}
  \item \emph{Say it:} \emph{pinchar}, once more
  \item \emph{Say it:} a sentence of your own with \emph{pinchar}
  \item \emph{Recall:} say the Spanish for the mechanic, then the Spanish for the tyre, then say \emph{pinchar} again
\end{itemize}

\begin{tcolorbox}[breakable,colback=teal!4,colframe=teal!35!black,title={Before you move on}]
What does pinchar mean? (\textquotedblleft{}To get a puncture; to prick; to click\textquotedblright{}.) Say it once more.
\end{tcolorbox}
